\chapter{RC and RL Circuits: Time Constants and AC Impedance}
\label{ch:rc-rl}

\section{The RC Charging Transient}

Connect a resistor and capacitor in series with a DC source (a step
voltage $U_0$ applied at $t=0$, capacitor initially uncharged). Solving
$i = C\,dU/dt$ together with Kirchhoff's voltage law gives the classic
exponential charging curve:
\[
\boxed{U_C(t) = U_0\left(1 - e^{-t/\tau}\right)}, \qquad \tau = RC
\]
where $\tau$ (tau), the \textbf{time constant}, has units of seconds when
$R$ is in ohms and $C$ is in farads. The current, meanwhile, starts at
its maximum value and decays exponentially:
\[
i(t) = \frac{U_0}{R}\,e^{-t/\tau}.
\]

\begin{center}
\begin{tikzpicture}
\begin{axis}[width=9.5cm,height=5cm,xlabel={$t/\tau$},ylabel={$U_C/U_0$},
  xmin=0,xmax=5,ymin=0,ymax=1.05,domain=0:5,samples=80,
  axis lines=left, ytick={0,0.632,1}, yticklabels={0,0.632,1.0}]
\addplot[thick,blue] {1-exp(-x)};
\draw[dashed] (axis cs:1,0) -- (axis cs:1,0.632) -- (axis cs:0,0.632);
\end{axis}
\end{tikzpicture}
\end{center}

\begin{keyidea}[Meaning of the time constant]
After one time constant ($t=\tau$), the capacitor has charged to
$1-e^{-1} \approx 63.2\%$ of the final voltage. After five time constants
($t=5\tau$), it is charged to $>99\%$ and is conventionally considered
``fully charged'' for practical purposes. The same 63.2\%/5-time-constant
rule applies to discharge (in reverse) and to the analogous RL circuit
below.
\end{keyidea}

\begin{examplebox}[RC time constant]
A 10\,k$\Omega$ resistor and a 100\,$\mu$F capacitor are placed in
series across a 9\,V source. Find the time constant and the voltage
across the capacitor after 2 seconds.

\emph{Solution.} $\tau = RC = 10{,}000 \times 100\times10^{-6} = 1\,\text{s}$.
At $t=2\,\text{s} = 2\tau$:
\[
U_C = 9\left(1-e^{-2}\right) = 9 \times 0.865 \approx 7.79\,\text{V}.
\]
\end{examplebox}

\section{RC Discharge}

If a charged capacitor is instead connected across a resistor (source
removed), the voltage decays exponentially from its initial value
$U_0$:
\[
U_C(t) = U_0\,e^{-t/\tau}, \qquad \tau = RC.
\]
This decaying-exponential shape, and its rising-exponential counterpart
above, appear throughout electronics: power-supply filter ripple decay,
the discharge of a defibrillator-style pulse capacitor, the decay of a
keyed CW envelope shaped by an RC network, and the timing element of many
oscillator and timer circuits.

\section{The RL Transient}

An inductor and resistor in series show the same mathematical shape,
because the inductor's $u=L\,di/dt$ plays the same differential-equation
role as the capacitor's $i=C\,dU/dt$, with current and voltage roles
exchanged. For current building up through an inductor after a voltage
step $U_0$ is applied:
\[
i(t) = \frac{U_0}{R}\left(1 - e^{-t/\tau}\right), \qquad \tau = \frac{L}{R}
\]
where now $\tau = L/R$ (seconds, when $L$ is in henries and $R$ is in
ohms). Current rises toward its final DC value ($U_0/R$, since the
inductor eventually looks like a plain wire) following the identical
63.2\%-per-time-constant pattern as the RC charging curve.

\bookfig[0.5]{rl_circuit.png}{A series RL circuit driven by a step voltage.}

\begin{examplebox}[RL time constant]
A relay coil has $L = 200\,\text{mH}$ and winding resistance
$R = 50\,\Omega$. Find the time constant, and the time for current to
reach 63.2\% of its final value.

\emph{Solution.} $\tau = L/R = 0.2/50 = 4\,\text{ms}$. By definition, the
current reaches 63.2\% of its final value at exactly $t = \tau =
4\,\text{ms}$.
\end{examplebox}

\section{Series RC and RL Circuits Under Sinusoidal AC}

Under a sinusoidal (not step) drive, a series RC or RL network behaves as
a frequency-dependent voltage divider between a real part (resistance)
and an imaginary part (reactance). Using the complex impedances
introduced in Chapters~\ref{ch:capacitors}--\ref{ch:inductors}, the
total impedance of a series RC network is
\[
Z = R - jX_C, \qquad |Z| = \sqrt{R^2 + X_C^2}, \qquad
\theta = -\arctan\!\left(\frac{X_C}{R}\right),
\]
and of a series RL network,
\[
Z = R + jX_L, \qquad |Z| = \sqrt{R^2 + X_L^2}, \qquad
\theta = \arctan\!\left(\frac{X_L}{R}\right).
\]
Here $\theta$ is the phase angle by which the total current lags (RL) or
leads (RC) the applied voltage. Ohm's law generalizes directly to
$U = IZ$, with $U$, $I$, and $Z$ all treated as complex (magnitude and
phase) quantities.

\begin{examplebox}[Series RC impedance]
A 1\,k$\Omega$ resistor is in series with a capacitor whose reactance at
the operating frequency is 750\,$\Omega$. Find the magnitude and phase
angle of the total impedance.

\emph{Solution.}
\[
|Z| = \sqrt{1000^2 + 750^2} = \sqrt{1{,}000{,}000+562{,}500} \approx 1250\,\Omega,
\]
\[
\theta = -\arctan\!\left(\frac{750}{1000}\right) = -36.9\degree.
\]
The current leads the applied voltage by $36.9\degree$.
\end{examplebox}

\section{RC and RL Circuits as Frequency-Selective Filters}

Because $X_C$ and $X_L$ depend on frequency while $R$ does not, a series
RC or RL network divides voltage between its resistive and reactive
parts differently at different frequencies --- the basis of simple
\textbf{filters}:

\begin{itemize}
\item An RC network, output taken across $C$, passes low frequencies and
attenuates high frequencies (a \textbf{low-pass filter}): at low
frequency $X_C$ is large and most of the voltage appears across $C$; at
high frequency $X_C \to 0$ and the output collapses.
\item An RC network, output taken across $R$, does the opposite (a
\textbf{high-pass filter}).
\item An RL network, output taken across $L$, is a high-pass filter (at
high frequency $X_L$ dominates); output taken across $R$ is a low-pass
filter.
\end{itemize}

The frequency at which $X_C = R$ (RC) or $X_L = R$ (RL) is called the
\textbf{cutoff (or corner) frequency} $f_c$, where the output power has
fallen to half its low-frequency (or high-frequency) value, a
$-3\,\text{dB}$ point in the language of Chapter~\ref{ch:decibels}:
\[
f_c = \frac{1}{2\pi RC} \quad \text{(RC)}, \qquad f_c = \frac{R}{2\pi L} \quad \text{(RL)}.
\]
These single-order RC/RL filters are gentle (attenuating at roughly
6\,dB per octave beyond cutoff); sharper filtering requires combining
reactive elements, covered next in Chapter~\ref{ch:rlc-resonance}.

\begin{quickreview}
\item RC time constant $\tau = RC$; RL time constant $\tau = L/R$; both
in seconds when SI units are used.
\item Charging/current-buildup: $63.2\%$ of final value after one time
constant; effectively complete after five time constants.
\item Series impedance: $Z = R \mp jX$ (RC/RL), $|Z|=\sqrt{R^2+X^2}$,
phase $\theta = \arctan(X/R)$ (sign per RC or RL).
\item RC/RL networks act as simple low-pass or high-pass filters
depending on which element the output is taken across.
\item Cutoff frequency $f_c$ is where $X = R$; the response is down
$3\,\text{dB}$ (half power) at that point.
\end{quickreview}

\begin{practiceproblems}
\item An RC circuit has $R=4.7\,\text{k}\Omega$ and $C=10\,\mu\text{F}$.
Find $\tau$ and the time to reach 99\% of final charge.
\item Find the impedance magnitude and phase of a series RL circuit with
$R=220\,\Omega$ and $X_L=330\,\Omega$.
\item Find the cutoff frequency of a low-pass RC filter with
$R=1\,\text{k}\Omega$ and $C=470\,\text{pF}$.
\item A relay must fully release within 20\,ms of power removal. If its
coil resistance is 100\,$\Omega$, what is the maximum inductance
(assuming ``fully released'' means five time constants have elapsed)?
\item Explain, without equations, why an RC low-pass filter's output
lags further and further behind the input as frequency increases.
\end{practiceproblems}
